\chapter{Open-Source Contributions as Proof-of-Work}
\label{ch:open-source}

Chapters~\ref{ch:projects}--\ref{ch:projects-c} built one form of proof: a production-grade project you own
end to end. This chapter builds a second, complementary form: proof that you can work inside \emph{someone
else's} codebase, collaborate with maintainers, and ship real, merged changes. The source material argues these
two forms of evidence are stronger together than either alone, and that both matter more than a resume for a
candidate without conventional employment history.

\section{Why Open Source, Specifically}

The starting problem: junior roles increasingly ask for a median of around two years of experience, and AI
coding tools let experienced engineers absorb work that used to go to juniors. A resume claiming ``Python,
machine learning, APIs, LLMs'' is easy to write and hard to verify. The proposed fix is to replace unsupported
claims with publicly inspectable evidence: merged contributions to established repositories, a deployed
application, tests and documentation, and a live demonstration of how you think.

\begin{caution}
The \$45,000--\$60,000/year (roughly INR 40--50 lakh) figure cited alongside this strategy is the
source's own estimate for remote AI-engineering compensation, not a verified market-wide range or a guarantee.
Actual pay depends on the company, experience, contract type, and location policy --- treat it the same way
Chapter~\ref{ch:foundations} treats its own cited salary range.
\end{caution}

\section{Choosing the Right Repository: The Repeatable-Unit Principle}

Generic ``contribute to open source'' advice fails for beginners because a huge repository with thousands of
open issues and no clear process produces weeks of confusion. The fix is to look for projects built from a
\textbf{repeatable unit} --- a small component you can implement by following an established pattern. Examples:
an LLM evaluation framework with dozens of independent metrics; an LLM API gateway supporting many providers
through a similar interface each time; a structured-output library with well-scoped modules and tests. Once you
understand one unit, the next one follows the same shape --- this is what makes a second and third contribution
dramatically faster than the first.

\subsection{Repositories Cited as Examples}

\begin{longtable}{p{0.18\textwidth} p{0.22\textwidth} p{0.52\textwidth}}
\toprule
\textbf{Repository} & \textbf{Focus} & \textbf{Why it was suggested} \\
\midrule
\endhead
DeepEval & LLM evaluation and testing & A metrics directory of roughly 52 independent metrics, each in its own
folder (e.g.\ answer relevancy, agent-loop detection) --- pick one, understand it, improve or extend it, write
tests. \\
LiteLLM & LLM API integrations & Roughly 138 provider integrations following a recurring pattern --- learn one,
use its structure to implement or improve another. \\
Instructor & Structured LLM outputs & Recommended for its focused scope and broad developer use --- look for a
reproducible bug, a missing example, or a narrowly scoped feature. \\
Agenta & Full-stack LLM tooling & A FastAPI backend plus React frontend --- useful if you want to demonstrate
more than Python-only skills. \\
DSPy & LLM programming frameworks & Cited for particularly clean documentation --- good for learning from a
well-organized codebase. \\
\bottomrule
\end{longtable}

\begin{caution}
These are the source's examples, not a required list, and its star counts and descriptions are not independently
verified here. Treat them as a starting point for applying the criteria below, not a checklist to complete.
\end{caution}

\subsection{Four Criteria for Choosing a Repository}

\begin{longtable}{p{0.26\textwidth} p{0.3\textwidth} p{0.36\textwidth}}
\toprule
\textbf{Criterion} & \textbf{What to look for} & \textbf{Why it matters} \\
\midrule
\endhead
Recent activity & Commits within the last week & Indicates active development \\
Genuine beginner issues & A \texttt{good first issue} label maintainers still actually use & Suggests the
project is open to new contributors \\
Contribution documentation & A real \texttt{CONTRIBUTING.md} & Explains how to submit acceptable changes \\
Manageable maintainer workload & A community that can realistically review your work & Improves the odds of
feedback and a timely merge \\
\bottomrule
\end{longtable}

As a rough heuristic on size: below $\sim$30{,}000 stars, maintainers tend to be more accessible; around
$\sim$56{,}000 stars (LiteLLM is the cited example) is a reasonable middle ground where the pattern is clear but
reviews take longer; above $\sim$100{,}000 stars (Dify, Open WebUI cited as examples), pull requests may sit in a
large queue. These are rules of thumb, not statistical laws --- maintainer responsiveness depends on the project
and the specific issue, not star count alone.

\section{The Exact Process for Getting Contributions Merged}

\begin{enumerate}
  \item \textbf{Read CONTRIBUTING.md first.} Understand the dev setup, coding conventions, test requirements,
  branch rules, and PR format. Many rejected PRs fail because contributors skip this, even when the code itself
  is reasonable.
  \item \textbf{Join the community before submitting code.} Discord or Slack if available --- introduce
  yourself, ask focused questions, understand maintainer expectations, build a real connection before opening a
  PR.
  \item \textbf{Make the first contribution small and boring.} Fix documentation, add a missing example, address
  a small bug with a clear reproduction. Don't attempt a massive feature to impress maintainers on your first
  try --- the goal is to learn the workflow and demonstrate reliability.
  \item \textbf{Repeat the same contribution pattern.} After your first change is accepted, implement another
  unit of the same type (another provider integration, another metric). This shows your first success wasn't a
  one-off.
  \item \textbf{Aim for three merged PRs in one recognizable repository.} Three accepted contributions to one
  project show sustained involvement and codebase familiarity, rather than scattered participation across
  unrelated projects.
\end{enumerate}

\begin{keytakeaway}
Plan for several weeks to get the \emph{first} PR merged --- you're learning the repo, the dev environment,
review conventions, and maintainer expectations all at once. Later contributions are typically much faster. The
important distinction is between \emph{submitting} a PR and \emph{getting it merged}: a submitted PR is not
proof the project accepted your work. Merged contributions, especially ones with substantive code or testing
changes, are the stronger evidence.
\end{keytakeaway}

\subsection{A Critical Warning About AI-Generated Contributions}

Do not copy a GitHub issue into an LLM, accept the generated solution, and submit it without understanding the
result. The risks are concrete: the proposed solution may not follow the project's architecture or conventions;
the code may be incorrect, poorly tested, or needlessly complicated; you may be unable to explain your own
implementation during review; maintainers may reject the contribution or lose confidence in you entirely.

\begin{caution}
This is not a caution about using AI --- it's a caution about using it to \emph{avoid} understanding your own
contribution. Several major open-source projects (cited examples: Ghostty, Godot) now have explicit public
policies on AI-generated PRs. Use AI instead to: explain unfamiliar modules and classes; help you understand
existing tests and architecture; interpret an error or help design test cases. Then write and review the
implementation yourself, run the relevant tests, and make sure you can explain and defend every important design
decision. The goal is not to avoid AI — it's to demonstrate engineering judgment while using it.
\end{caution}

\section{Build One Complete Product, Too}

Open-source contributions prove you can work inside someone else's project. You should also prove you can build
something independently end to end — frontend, backend, a data/model layer where relevant, testing, real
deployment, and ideally real users (the detailed how-to for this is Chapters~\ref{ch:projects}--\ref{ch:projects-c}).
The reasoning cited for why founders specifically value this: AI coding tools make producing code less
distinguishing, so what founders actually need to know is whether you can identify the right problem, make
sensible trade-offs, connect components into a working system, debug things with no obvious solution, and ship
without constant supervision. That combination — not raw code volume — is what ``judgment and ownership'' means
in practice.

\section{Open Source as a Networking Strategy}

One claim cited in the source material: roughly 58\% of hires across a set of AI-startup founding-engineer
searches came through founders' prior networks.

\begin{caution}
The source does not provide the underlying study, sample size, or methodology — treat this as a claim from one
video, not a verified statistic for the AI job market generally.
\end{caution}

The mechanism proposed is still useful regardless of the exact percentage: open-source work creates visibility
to maintainers (a merged contribution gives an experienced engineer a reason to recognize your name); it puts
you inside relevant communities (Discord/Slack servers may include people working at AI startups); it creates a
public artifact a founder can inspect instead of taking your resume at face value; and it builds a relationship
\emph{before} you ask for anything. Networking, on this view, is not a separate activity that starts once your
portfolio is finished — it emerges directly from the work itself.

\section{Overcoming the Three Objections to Hiring Remotely from India}

Even with the right skills, a founder may hesitate to hire remotely. Three specific objections, each with a
concrete counter:

\subsection{Objection 1: Time-Zone Overlap}

India runs roughly 9.5--13.5 hours ahead of the US, depending on coast and daylight saving. Many companies claim
to be asynchronous but still rely on synchronous collaboration in practice.

\begin{template}
I work from 2\,PM to 11\,PM IST, providing a consistent four-hour overlap with your team's morning.
\end{template}

Calculate the actual overlap for the specific company and season rather than copying this line blindly, and only
promise hours you can genuinely sustain — demonstrate the commitment by actually showing up during that window
throughout the hiring process.

\subsection{Objection 2: Founders Can't Easily Assess Remote Candidates}

Founders traditionally judge candidates through live, unscripted interviews — watching how someone reasons under
uncertainty. In an environment where AI-assisted coding and interview assistance are common, that trust problem
is sharper for remote candidates.

\begin{template}
I'd be happy to screen-share for 45 minutes and build something in front of you, without extensive preparation.
\end{template}

A strong session involves reading unfamiliar code, asking clarifying questions, designing a small feature,
implementing it, writing a test, and explaining the trade-offs out loud. Agree on a reasonable scope up front,
and don't let this turn into open-ended unpaid production work.

\subsection{Objection 3: Hiring Risk and Uncertainty About Quality}

A cheap hire becomes expensive fast if they need excessive supervision or produce code other engineers must
rewrite. The founder's decision is about expected cost and risk, not just salary. The recommended counter is to
combine all three pieces of evidence this chapter and Part~III build: three merged PRs in a recognizable
repository, one genuinely deployed end-to-end application, and the willingness to demonstrate live. Each piece
addresses a different uncertainty — collaboration ability, ownership, and real-time thinking, respectively. The
underlying strategy: don't argue you're a good candidate — give the employer evidence that lets them reach that
conclusion themselves.

\section{Cold Emails, LinkedIn, and Job Boards — This Source's Numbers}

The source quotes an average cold-email reply rate of roughly 3.4\% (as of 2026), lower than a cited historical
figure, with LinkedIn messaging said to perform better without a clear comparable dataset.

\begin{caution}
Treat 3.4\% as directional, not predictive for you individually — response rates vary heavily with targeting,
message quality, sender reputation, and whether the recipient is actually hiring. At a 3.4\% reply rate, emailing
1,000 companies produces roughly 34 replies on average — not 34 interviews or offers, and likely fewer replies
still for a candidate with no publicly verifiable work. This is one of three distinct quantitative models for
outreach in this guide; see the footnote in Chapter~\ref{ch:checklists} for how it relates to the other two.
\end{caution}

The recommended sequence mirrors Chapter~\ref{ch:direct-outreach}'s structure but is tuned for this strategy
specifically: identify a startup whose product and engineering needs match your skills $\rightarrow$ find a
relevant open-source repository, maintainer, or engineer in that ecosystem $\rightarrow$ make your work easy to
inspect by linking merged contributions and your deployed app $\rightarrow$ send a concise, personalized message
explaining why you're contacting this person and what you've built $\rightarrow$ state your working-hour overlap
and offer a live technical session if useful $\rightarrow$ follow up a limited number of times (three total
touches is cited as capturing most of the replies an outreach sequence will get), then stop.

Job boards specifically recommended alongside this strategy: \textbf{Work at a Startup}
(\url{workatastartup.com}, YC's own board) and \textbf{Wellfound} (\url{wellfound.com}, formerly AngelList
Talent) — used alongside public engineering work and targeted networking, not as a replacement for either.

\section{The Six-Phase Roadmap}

This roadmap runs in parallel to, not instead of, the 16-week roadmap in Chapter~\ref{ch:roadmap} — the two
describe the same underlying effort (build proof, get visible, reach out) through slightly different lenses
(project-first vs.\ open-source-first).

\begin{longtable}{p{0.08\textwidth} p{0.24\textwidth} p{0.6\textwidth}}
\toprule
\textbf{Phase} & \textbf{Timing} & \textbf{Focus} \\
\midrule
\endhead
1 & Days 1--3 & Select your ecosystem: choose an AI domain (LLM apps, evaluation, API integrations); identify
one active repo with a clear process; read its CONTRIBUTING.md and inspect recent merged PRs. \\
2 & Week 1 onward & Make your first contribution: set up the repo and run its tests; join the community; choose
a small, reproducible issue; submit and respond professionally to review feedback. \\
3 & Following weeks & Establish a contribution record: get the first PR merged, then repeat the pattern; work
toward three substantive merged PRs in the same repo; document what changed, why, and what you learned. \\
4 & In parallel & Build your own product (Chapters~\ref{ch:projects}--\ref{ch:projects-c}): choose a problem a
specific user type actually has; build, test, deploy; gather real feedback; prepare a short demo and a clear
explanation of your architecture. \\
5 & Once you have evidence & Become discoverable: update GitHub and LinkedIn with your strongest work; write a
post about what you learned; build genuine relationships with maintainers; prepare a concise portfolio page
linking PRs, code, tests, and the live app. \\
6 & Ongoing & Apply and demonstrate: target relevant startups via careers pages, Work at a Startup, and
Wellfound; send personalized outreach referencing the company's needs; state your real working-hour overlap;
offer a live session; review the contract, payment arrangement, and tax obligations before accepting (see
Chapter~\ref{ch:getting-paid}). \\
\bottomrule
\end{longtable}

\section{What to Trust, and What to Treat Cautiously}

\begin{longtable}{p{0.44\textwidth} p{0.5\textwidth}}
\toprule
\textbf{Claim or recommendation} & \textbf{How to interpret it} \\
\midrule
\endhead
Remote AI roles can pay \$45,000--\$60,000 & A possible range, not a universal entry-level salary — verify
individual roles and your own eligibility \\
Three merged PRs make you a strong candidate & A useful initial portfolio target, not a hiring threshold —
contribution quality and relevance matter more than the raw count \\
Smaller repos provide better maintainer access & A reasonable heuristic, but community culture and reviewer
activity are more informative than star count alone \\
58\% of hires come through founders' networks & A claim about the searches described in the source material; it
does not establish how representative this is more broadly \\
Three outreach touches capture almost all replies & Not a universal recruiting rule — follow-up effectiveness
depends on context, and repeated unsolicited messages can backfire \\
Contractors can avoid US withholding using W-8BEN & Incomplete without examining the type of income, treaty
eligibility, and documentation — see Chapter~\ref{ch:getting-paid} \\
\bottomrule
\end{longtable}

\begin{keytakeaway}
The entire strategy, compressed: pick one relevant, active open-source project; learn its architecture and
contribution process; get three meaningful PRs merged, ideally in the same repo; build and deploy one
end-to-end product; publish your work and build genuine relationships with engineers and maintainers; apply
selectively, demonstrate live, and state your working-hour overlap clearly; understand the legal, payment, and
tax implications before accepting an offer. The most important distinction in this whole chapter is between
\emph{claiming} you can do the job and making it \emph{possible for a founder to verify} that you can.
\end{keytakeaway}

